\pg{336}

\textbf{60. O ente contingente é produzido, immediata ou mediatamente, pelo ente necessario.}

\textbf{a)} A causa extrinseca, que produz o ente \emph{contingente}, é um

ente \emph{necessario}; ou \emph{contingente}. Se é \emph{necessario}, está provada a conclusão. Se é \emph{contingente}, este foi produzido por um ente \emph{necessario}, ou por um ente \emph{contingente}. Se foi produzido por um ente \emph{necessario}, temos o que queriamos. Se foi produzido por um ente \emph{contingente}, renovamos a pergunta, e, subindo de um ente para outro, devemos inevitavelmente encontrar o ente \emph{necessario}; porque, se não existisse o ente \emph{necessario}, não poderia existir nenhum ente \emph{contingente}, que não tem em si a razão da sua existencia. Logo o ente \emph{contingente} é produzido, immediata ou mediatamente, pelo ente \emph{necessario} \footnote{A conclusão é um corollario da precedente. O ente contingente, por ser \emph{contingente}, não tem em si a razão da propria existencia; logo deve tel-a numa causa extrinseca, no ente \emph{necessario}. --- Dizemos que o ente \emph{contingente} depende, \emph{immediata} ou \emph{mediatamente}, do ente \emph{necessario}; porque, embora um ente \emph{contingente} possa ser produzido immediatamente por um outro ente \emph{contingente}, comtudo a ultima razão adequada da existencia do mesmo ente \emph{contingente} não pode encontrar-se senão no ente \emph{necessario}.}.

\textbf{b)} O ente, que se acha em potencia em relação a um certo acto, não o pode adquirir, se não lh'o fôr communicado, immediata ou mediatamente, por outro ente, que o possua pela sua essencia. Ora o ente \emph{contingente}, sendo por si indifferente para a existencia e para a não existencia, não existiu sempre e não pode adquirir a existencia, se não lh'a fôr communicada, immediata ou mediatamente, pelo ente \emph{necessario}, que a possua pela exigencia da sua essencia. Logo o ente \emph{contingente} é produzido, immediata ou mediatamente, pelo ente \emph{necessario} \footnote{O ente \emph{contingente} é um ente \emph{em potencia}, isto é, \emph{pode} receber o acto da existencia, mas não o possue por si, pela exigencia da sua natureza. Ora um ente, que não possue por si uma perfeição, não a pode communicar a si mesmo, isto é, não pode passar por si da potencia para o acto. Logo essa perfeição deve ser-lhe communicada por outro ente. Este ente não poderia\pgn{337} communicar a outro uma perfeição, se a não possuisse, porque ninguem dá o que não tem. Porquanto, o ente \emph{contingente}, que é indifferente para a existencia e a não-existencia, que está em potencia em relação á existencia, ficaria sempre nessa indifferença, nunca poderia adquirir a existencia, se não lhe fosse communicada pelo ente que a possue pela sua essencia, isto é, pelo ente \emph{necessario}. (Cf. \emph{C. Gent}, I, 16). --- Alem d'isso, o ente \emph{contingente} é ente \emph{por participação}, ou \emph{participado}; porque não possue o ser pela propria essencia, não é o proprio ser. Ora o ente \emph{participado} deve derivar do ente, que é tal pela sua essencia; isto é, do ente \emph{imparticipado}; porque o ente \emph{participado}, não sendo o proprio ser, recebeu o ser de outro ente, que, em ultima analyse, deve possuir o ser pela propria essencia. Diz S. Thomaz: « Necesse est dicere omne id quod praeter \textsc{Deum} est, a \textsc{Deo} factum esse. Solus enim Deus est suum \emph{esse}; in omnibus autem aliis differt essentia rei et esse ejus... Et ex hoc manifestum est quod solus Deus est ens per suam essentiam; omnia vero alia sunt entia per participationem. Omne autem quod est per participationem, causatur ab eo quod est per essentiam; sicut omne ignitum causatur ab igne » (\emph{Sum. Th.}, p. I, q. 61, a. 1; Cf. \emph{C. Gent.}, I, 15).}.
